\subsection{ASSOCIATIVITY OF THE TENSOR PRODUCT}

PROPOSITION 8. Let A, B be two rings, E a right A-module, F an (A, B)-bimodule and G a left B-module. Then E \(\otimes_{\mathbf{A}} \mathbf{F}\) is a right B-module, F \(\otimes_{\mathbf{B}} \mathbf{G}\) a left A-module and there exists one and only one Z-linear mapping

\[
\phi : (\mathrm{E} \otimes_ {\mathrm{A}} \mathrm{F}) \otimes_ {\mathrm{B}} \mathrm{G} \rightarrow \mathrm{E} \otimes_ {\mathrm{A}} (\mathrm{F} \otimes_ {\mathrm{B}} \mathrm{G})
\]

such that \(\phi((x \otimes y) \otimes z) = x \otimes (y \otimes z)\) for \(x \in \mathbf{E}\) , \(y \in \mathbf{F}\) , \(z \in \mathbf{G}\) ; moreover this \(\mathbf{Z}\) -linear mapping is bijective (``associativity'' of the tensor product).

The right B-module structure on \(\mathbf{E} \otimes_{\mathbf{A}} \mathbf{F}\) and left A-module structure on \(\mathbf{F} \otimes_{\mathbf{B}} \mathbf{G}\) have been defined in no. 4. The uniqueness of \(\phi\) is obvious since the elements \((x \otimes y) \otimes z\) generate the Z-module \((\mathbf{E} \otimes_{\mathbf{A}} \mathbf{F}) \otimes_{\mathbf{B}} \mathbf{G}\) . To show the existence of \(\phi\) , we note that, for all \(z \in \mathbf{G}\) , \(h_z: y \mapsto y \otimes z\) is an A-linear mapping of the left A-module F into the left A-module \(\mathbf{F} \otimes_{\mathbf{B}} \mathbf{G}\) . We write \(g_z = l_{\mathbf{E}} \otimes h_z\) , which is therefore a Z-linear mapping of \(\mathbf{E} \otimes_{\mathbf{A}} \mathbf{F}\) into \(\mathbf{E} \otimes_{\mathbf{A}} (\mathbf{F} \otimes_{\mathbf{B}} \mathbf{G})\) and consider the mapping \((t, z) \mapsto g_z(t)\) from \((\mathbf{E} \otimes_{\mathbf{A}} \mathbf{F} \times \mathbf{G}\) into \(\mathbf{E} \otimes_{\mathbf{A}} (\mathbf{F} \otimes_{\mathbf{B}} \mathbf{G})\) ; as \(h_{z+z'} = h_z + h_{z'}\) for \(z \in \mathbf{G}\) , \(z' \in \mathbf{G}\) , it is immediate that the above mapping is Z-bilinear. Further, we show that for all \(\mu \in \mathbf{B}\) , \(g_{\mu z}(t) = g_z(t\mu)\) ; it is obviously sufficient to do this for \(t = x \otimes y\) where \(x \in \mathbf{E}\) and \(y \in \mathbf{F}\) ; now

\[
g _ {\mu z} (x \otimes y) = x \otimes (y \otimes \mu z)
\]

\[
g _ {z} ((x \otimes y) \mu) = g _ {z} (x \otimes y \mu) = x \otimes (y \mu \otimes z).
\]

Proposition 1 (no. 1) then proves the existence of a Z-linear mapping

\[
\phi : (\mathrm{E} \otimes_ {\mathrm{A}} \mathrm{F}) \otimes_ {\mathrm{B}} \mathrm{G} \rightarrow \mathrm{E} \otimes_ {\mathrm{A}} (\mathrm{F} \otimes_ {\mathrm{B}} \mathrm{G})
\]

such that \(\phi(t \otimes z) = g_z(t)\) , hence \(\phi((x \otimes y) \otimes z) = x \otimes (y \otimes z)\) . Similarly a Z-linear mapping

\[
\psi : \mathbf {E} \otimes_ {\mathbf {A}} (\mathbf {F} \otimes_ {\mathbf {B}} \mathbf {G}) \rightarrow (\mathbf {E} \otimes_ {\mathbf {A}} \mathbf {F}) \otimes_ {\mathbf {B}} \mathbf {G}
\]

is defined such that \(\psi(x \otimes (y \otimes z)) = (x \otimes y) \otimes z\) and clearly \(\psi \circ \phi\) and \(\phi \circ \psi\) are the identity mappings of \((\mathbf{E} \otimes_{\mathbf{A}} \mathbf{F}) \otimes_{\mathbf{B}} \mathbf{G}\) and \(\mathbf{E} \otimes_{\mathbf{A}} (\mathbf{F} \otimes_{\mathbf{B}} \mathbf{G})\) respectively, since they reduce to the identity mappings on generating systems of these \(\mathbf{Z}\) -modules.

It is immediate that, if E is a \(((C_{i}^{\prime}); A, (D_{j}^{\prime}))\) -multimodule, F an \((A, (C_{h}^{\prime\prime}); B, (D_{k}^{\prime\prime}))\) -multimodule and G a \((B, (C_{l}^{\prime\prime}); (D_{m}^{\prime\prime}))\) -multimodule, the canonical isomorphism defined in Proposition 8 is a \(((C_{i}^{\prime}), (C_{h}^{\prime\prime}), (C_{l}^{\prime\prime}); (D_{j}^{\prime}), (D_{k}^{\prime\prime}), (D_{m}^{\prime\prime}))\) -multimodule isomorphism. In particular, if C is a commutative ring and E, F, G three C-modules, there is a canonical C-module isomorphism

\[
(\mathbf {E} \otimes_ {\mathbf {c}} \mathbf {F}) \otimes_ {\mathbf {c}} \mathbf {G} \rightarrow \mathbf {E} \otimes_ {\mathbf {c}} (\mathbf {F} \otimes_ {\mathbf {c}} \mathbf {G}).
\]

We shall see below that, under certain conditions, the definition of tensor product can be generalized to a family of multimodules, which will in particular give us under the hypotheses of Proposition 8 a Z-module \(E \otimes_{A} F \otimes_{B} G\) , which is canonically isomorphic to each of the Z-modules \((E \otimes_{A} F) \otimes_{B} G\) and \(E \otimes_{A} (F \otimes_{B} G)\) and with which the latter will be identified.

